\chapter{Data Analysis Script}\label{appendix:data_analysis_script}
\addcontentsline{toc}{chapter}{APPENDIX~\thechapter: Data Analysis Script}

The following Python script was used to analyze the stress-relaxation data and generate the plots presented in this dissertation. The script processes CSV files containing force-displacement data, performs curve fitting to extract mechanical parameters, and outputs results to Excel files.

\subsection{Key Features}
\begin{itemize}
    \item Automated processing of multiple CSV files
    \item Exponential curve fitting for stress-relaxation analysis 
    \item Calculation of key mechanical parameters (max force, time constants, moduli)
    \item Generation of force vs. time plots with curve fits
    \item FFT analysis of force signals
    \item Export of analyzed data to Excel spreadsheets
\end{itemize}

\subsection{Script Implementation}

\subsubsection{Required Packages}
\begin{lstlisting}[caption={Required package imports}]
import pandas as pd
import glob
import matplotlib.pyplot as plt
import numpy as np
import sympy as sym
from scipy.optimize import curve_fit
import openpyxl
from scipy.fft import fft, fftfreq
import easygui as gui
import sys
\end{lstlisting}

\subsubsection{Curve Fitting Functions}
\begin{lstlisting}[caption={Curve fitting function definitions}]
# Exponential function for curve fitting
def exp(x, a, b, c):
    return a * np.exp(-x/b) + c

# Linearized exponential for regression
def linear_exp(x, a, b):
    return -(x/b) + np.log(a)
\end{lstlisting}

\subsubsection{Global Constants and Settings}
\begin{lstlisting}[caption={Global constants and settings}]
# Data column headers
SETNAME = 'SetName'
CYCLE = 'Cycle'
TIME_S = 'Time_S'
SIZE_MM = 'Size_mm' 
DISP_MM = 'Displacement_mm'
FORCE_N = 'Force_N'
FORCE_N_EMA = 'Force_N_EMA'
TIME_S_HOLD = 'Time_S_hold'

# Physical and experimental parameters
Area_m2 = 0.0001267  # Sample area in m^2
strain_map = {
    '5%' : 0.05,
    '10%' : 0.1,
    '20%' : 0.2,
}
strain_values = ['5%', '10%', '20%']

# Analysis parameters
SPAN_FORCE = 5  # Smoothing parameter
COMPRESS_REGION = '1-Compress'
HOLD_REGION = '1-Hold'
SAMPLE_RATE = 100
\end{lstlisting}

\subsubsection{File Handling and Setup}
\begin{lstlisting}[caption={File handling and Excel setup}]
# GUI folder selection
folderPath = gui.diropenbox("Choose folder with .csv files!")
if folderPath is None:
    sys.exit(0)

# Excel file handling
msg = "Do you want to create a new excel file or append to an existing excel file\nOr Press <cancel> to exit."
createChoice = gui.choicebox(msg, "Excel file choice", ["Create new", "Append"])

if createChoice is None:
    sys.exit(0)

if createChoice == "Append":
    excelFile = gui.fileopenbox(filetypes=['*.xlsx'], 
                               title="Select properly formatted existing excel file!")
else:
    newExcelFolder = gui.diropenbox("Choose folder to save new excel file!")
    newExcelName = gui.enterbox("Enter name of new excel file!", 
                               "Name input", "DataSummary")
    newExcelPath = newExcelFolder + '\\' + newExcelName + '.xlsx'
    writer = pd.ExcelWriter(newExcelPath, engine='auto')

    # Initialize Excel sheets for each strain value
    for workbook_name in strain_values:
        df = pd.DataFrame({
            'Name': [], 'Fmax': [], "~0.3679*Fmax": [],
            'Estimated time constant': [], 
            'a (Fmax - Fequil. for eq.)': [],
            'b (time constant for eq.)': [],
            'c (equilibrium force for eq.)': [],
            'Aggregate Modulus (HA)': [],
            'Area under hold start to Fmax/2 + Feq/2': [],
            'Area under entire hold region': [],
            'Impulse over half area': [],
            'Impulse over full area': []
        })
        df.to_excel(writer, sheet_name=workbook_name, index=False)

    writer.close()
    excelFile = newExcelPath
\end{lstlisting}

\subsubsection{Main Analysis Loop}
\begin{lstlisting}[caption={Main analysis implementation}]
# Process all CSV files
files = glob.glob(folderPath + '/*.csv')
wb = openpyxl.load_workbook(excelFile)

for filename in files:
    df_all = pd.read_csv(filename)
    for strain in strain_values:
        # Analysis implementation continues as in original script...
        # Code truncated for brevity
\end{lstlisting}
\newpage
\subsection{Usage Notes}
The script requires input CSV files containing stress-relaxation test data with specific column headers, including time, displacement, and force measurements. It processes multiple strain levels (5\%, 10\%, 20\%) and generates both graphical and numerical outputs. The graphical outputs include stress-relaxation curves and model fits, while numerical outputs comprise parameters like maximum force, time constants, and aggregate modulus. The analyzed data is automatically saved to Excel files for further processing, with separate worksheets for each strain level. The script provides an interactive interface for selecting input files and output locations through dialog boxes.